\chapter{Родитель внешнего квадрата обменного моста и аудит нормировки}

Предыдущая проверка вернула из Тома IV точный инвариант разности прямого и перекрёстного каналов
$4e_2(R)$ и показал, что на уровне состояния он создаёт одноосную
$\mathbb{RP}^2$-фазу. Теперь проверяется более сильное утверждение:
возникает ли этот член из уже выбранного минимального Real-родителя, и
сохраняется ли переход после добавления полного самосогласованного
мостового действия.

\section{Литературная и проектная граница}

Суперсвязность Квиллена допускает нечётную нуль-форменную эндоморфную
часть, а её кривизна является обычным градуированным квадратом; см.
Quillen, \emph{Superconnections and the Chern character}, Topology 24
(1985), DOI \path{10.1016/0040-9383(85)90047-3}. Индуцированная связность
на внешней алгебре является стандартной функториальной конструкцией; см.
\path{arXiv:hep-th/9801045}, \path{arXiv:math-ph/9911006} и
\path{arXiv:1308.3275}.

Однако эти источники различают два объекта:
\begin{enumerate}
 \item аддитивный подъём оператора на внешнюю алгебру;
 \item мультипликативную внешнюю степень $\Lambda^2B$ конечного перехода.
\end{enumerate}
Вторая квадратична по $B$ и не обязана входить в кривизну исходной
двухузловой суперсвязности.

Проектный прецедент также имеет точную границу. проверки Пати--Салама вывели
внешний квадрат только после добавления относительной трёхузловой цепи и
фактор-нормы по алгебре неподвижных точек. Обычный полный след оставался
ранг-слепым. Поэтому нельзя переносить итоговый коэффициент четыре без
повторной проверки носителя, следа и кратности.

\section{Обычная двухузловая суперсвязность не видит внешний квадрат}

На семейных половинах минимальный нечётный эндоморфизм имеет вид
\begin{equation}
 \Phi_B=
 \begin{pmatrix}0&B^T\\B&0\end{pmatrix}.
\end{equation}
Если $X=B^TB$, то
\begin{equation}
 \Phi_B^2=\operatorname{diag}(X,BB^T),
 \qquad
 \Tr\Phi_B^4=2\Tr X^2.
 \label{eq:v6-two-node-bridge-trace-four}
\end{equation}
Нужный инвариант равен
\begin{equation}
 e_2(X)=\frac12\bigl((\Tr X)^2-\Tr X^2\bigr).
\end{equation}
Таким образом, обычный односвязный след даёт $\Tr X^2$, но не создаёт
второй след $(\Tr X)^2$. Это тот же структурный разрыв между полным и
относительным действием, который уже встречался в Томах IV--V.

\begin{theorem}[Двухузловой запрет]
Кривизна минимальной Real-суперсвязности с одним мостом $B$ не содержит
$e_2(B^TB)$. Следовательно, найденный внешний квадрат не является
скрытым членом уже принятого канонического действия.
\end{theorem}

\section{Минимальный функториальный носитель внешней степени}

Пусть $U=V=\mathbb R^3$. Рассмотрим цепь
\begin{equation}
 \mathbb R
 \xrightarrow{\ A_B\ }
 U\otimes V
 \xrightarrow{\ C_B\ }
 \Lambda^2U\otimes\Lambda^2V,
 \label{eq:v6-bridge-exterior-chain}
\end{equation}
где
\begin{equation}
 A_B(1)=\operatorname{vec}B,
 \qquad
 C_B=\frac12d(\Lambda^2)_B.
\end{equation}
Множитель $1/2$ фиксирован поляризацией квадратичной карты:
\begin{equation}
 C_BA_B=\Lambda^2B.
 \label{eq:v6-polarized-exterior-path}
\end{equation}
Для стандартных ортонормированных базисов внешних степеней
\begin{equation}
 \|\Lambda^2B\|_F^2=e_2(B^TB).
 \label{eq:v6-normalized-exterior-norm}
\end{equation}

На самосопряжённом операторе цепи размерности $1+9+9$ чётная кривизна
содержит путь \eqref{eq:v6-polarized-exterior-path} и сопряжённый путь.
Относительная фактор-норма блоков конечных узлов поэтому равна
\begin{equation}
 \boxed{
 S_{\Lambda^2}^{\rm rel}(B)
 =2e_2(B^TB).}
 \label{eq:v6-canonical-exterior-relative-action}
\end{equation}
Трёхмерная особенность
$\Lambda^2\mathbb R^3\simeq\mathbb R^3$ позволяет записать конечный путь
как матрицу кофакторов. Но она не устраняет новый узел цепи и не превращает
квадратичную карту $B\mapsto\Lambda^2B$ в линейное представление
$M_3(\mathbb R)$.

После нормировки $R=X/\Tr X$ формула
\eqref{eq:v6-canonical-exterior-relative-action} даёт
\begin{equation}
 2e_2(R)=1-\Tr R^2.
 \label{eq:v6-canonical-exterior-state-weight}
\end{equation}
То есть эквивалентный коэффициент отрицательного инварианта чистоты равен
единице, а не двум.

\section{Полный фазовый порог}

Предыдущий зависящий только от состояния расчёт сравнивал энтропию с внешним квадратом. Для
родительского вывода нужно вернуть самосогласованную радиальную энергию.
Пусть $p_i$ --- собственные значения $R$, а $x_i=Tp_i$. Минимизация по
общей амплитуде даёт
\begin{equation}
 T_*(R)=\frac{\sum_i p_i^2}{\sum_i p_i^3},
 \qquad
 S_{\rm rad}(R)=\frac27\left(
 1-\frac{(\sum_i p_i^2)^2}{\sum_i p_i^3}
 \right).
 \label{eq:v6-self-consistent-radial-elimination}
\end{equation}
Полный однородный функционал с внешнеквадратным весом $m$ равен
\begin{equation}
 \mathcal F_m(R)
 =S_{\rm rad}(R)+\Tr(R\log R)+m e_2(R).
 \label{eq:v6-full-exterior-free-energy}
\end{equation}

На одноосной семье
\begin{equation}
 R(a)=\operatorname{diag}\left(a,\frac{1-a}{2},\frac{1-a}{2}\right)
\end{equation}
точный численная проверка сосуществования фаз даёт
\begin{equation}
 \boxed{m_{\rm crit}=3.0261454269\ldots}
 \label{eq:v6-full-exterior-threshold}
\end{equation}
при $a=0.880484\ldots$. Эквивалентный коэффициент чистоты равен
$m_{\rm crit}/2=1.5130727\ldots$.

Каноническая относительная цепь имеет $m=2<m_{\rm crit}$ и сохраняет
изотропный минимум. Поэтому одного функториального носителя внешней степени
недостаточно.

\section{Почему сырой прямой минус перекрёстный канал проходит}

Если не переходить к ортонормированному базис внешней степениу, а хранить полный
четырёхиндексный тензор
\begin{equation}
 W_{ia,jb}=B_{ia}B_{jb}-B_{ib}B_{ja},
\end{equation}
то обычная фробениусова норма даёт
\begin{equation}
 \|W\|_F^2=4e_2(B^TB).
 \label{eq:v6-raw-tensor-exterior-factor-four}
\end{equation}
После нормировки состоянием это соответствует $m=4>m_{\rm crit}$.
Полный функционал \eqref{eq:v6-full-exterior-free-energy} тогда имеет
одноосный минимум
\begin{equation}
 \Spec R_*
 =(0.9668900\ldots,0.0165550\ldots,0.0165550\ldots),
 \label{eq:v6-full-raw-wedge-spectrum}
\end{equation}
а оптимальная общая амплитуда равна
\begin{equation}
 T_*=1.0348398\ldots.
\end{equation}
Тем самым ненормированный канал действительно выдерживает полный, а не только
зависящий только от состояния фазовый тест.

Но множитель два между $2e_2$ и $4e_2$ пока не выведен. В
ортонормированной внешней алгебре он исчезает как двойной подсчёт
антисимметричных компонент. KO6-удвоение не возвращает его: физический
полуслед сокращает удвоение. Вторая одинаковая относительная копия также
запрещена прежним критерием неприводимости, если не добавить новую
динамику пространства копий.

\begin{nogocriterion}[Запрет метрического удвоения]
Нельзя выбирать ненормированную тензорную норму вместо нормированной нормы внешней степени только
для прохождения порога. Дополнительный множитель два должен следовать из
самостоятельного пути между конечными узлами, ориентационной пары, условного ожидания или
иного оператора текущего родителя. Иначе это скрытый вес сектора.
\end{nogocriterion}

\section{Вердикт}

\begin{protocol}
\begin{enumerate}
 \item внешний квадрат в обычной двухузловой кривизне:
 \textbf{отсутствует};
 \item функториальный носитель внешней степени: \textbf{построен};
 \item двухшаговый путь: \textbf{$\Lambda^2B$};
 \item каноническая относительная норма: \textbf{$2e_2$};
 \item полный порог перехода: \textbf{$m_{\rm crit}=3.026145\ldots$};
 \item достаточность канонической одной копии: \textbf{нет};
 \item ненормированная норма разности прямого и перекрёстного каналов: \textbf{$4e_2$};
 \item прохождение полного фазового теста ненормированным каналом: \textbf{да};
 \item происхождение дополнительного множителя два:
 \textbf{не доказано};
 \item точный ранг один: \textbf{не требуется};
 \item полноранговая $\mathbb{RP}^2$-фаза: \textbf{условно получена};
 \item текущий минимальный родитель: \textbf{недостаточен};
 \item следующая проверка:
 \textbf{нормировка относительного носителя тензорного квадрата};
 \item рождение материи: \textbf{не закрыто}.
\end{enumerate}
\end{protocol}

Следующая проверка должна выяснить, содержит ли тензорный квадрат двух
обменных ориентаций два независимых пути между конечными узлами, которые канонически
сохраняют ненормированный коэффициент четыре после Real-полуследа. Если нет,
внешнеквадратная ветвь остаётся эффективной моделью, а не выведенным
родителем.

Машинный сертификат сохранён в
\path{s2t_v6_exchange_bridge_exterior_square_parent_gate_results.json}.